\reading{American Ambassador Casualty Co. v. City of Chicago}
\readingcite{563 N.E.2d 882 (Ill. App. 1990)}

\opinion Presiding Justice \textsc{LaPorta} delivered the opinion of the
court\ldots

%The subject of this appeal is the dismissal with prejudice of plaintiff's
%complaint which sought recovery from defendant of \$10,552, the amount paid by
%plaintiff to its insured for the loss of the insured's automobile. The trial
%court granted defendant's motion to dismiss which asserted that defendant was
%immune from suit under the Local Governmental and Governmental Employees Tort
%Immunity Act (Tort Immunity Act) (Ill.Rev.Stat.1987, ch. 85, par. 4-102).
%Plaintiff appeals, asserting that (1) defendant was liable under a bailment
%theory, (2) the Tort Immunity Act was inapplicable, and (3) even if the Tort
%Immunity Act was applicable, defendant was liable for breach of its special
%duty
%owed to plaintiff's insured.

Plaintiff
%, American Ambassador Casualty Company, as subrogee of its insured,
Bruce Anderson\edfootnote{The actual plaintiff of this case was Anderson's
insurer. To simplify matters, the opinion has been edited to make Anderson the
plaintiff. The difference is irrelevant to the analysis.}
%,
brought this action against defendant, the City of Chicago,
seeking recovery of \$10,552.50, the [insurance value]
%amount paid to Anderson for the loss
of his
vehicle. The complaint alleged that\ldots
% Anderson's vehicle was insured by plaintiff
%against loss by theft under an automobile policy providing comprehensive
%coverage. Plaintiff asserted further that
defendant impounded Anderson's
automobile after the driver of the vehicle, not Anderson, was arrested on
December 22, 1988, and although Anderson demanded the return of his vehicle,
defendant was unable or unwilling to return it. The automobile was reported
stolen on January 26, 1989, by Chicago police officer [Raymond]
Pellegrini.\ldots
%Plaintiff sought a judgment from defendant for the \$10,052 paid to Anderson,
%plus the \$500 deductible under the insurance policy.

%Defendant filed a motion to dismiss the complaint, asserting that it was immune
%from suit under the Tort Immunity Act (Ill.Rev.Stat.1987, ch. 85, par. 4-102).
%Plaintiff responded to defendant's motion, asserting that defendant was liable
%under a bailment theory or, in the alternative, for breach of the special duty
%owed to plaintiff's insured. Defendant replied, arguing that no bailment
%relationship had been shown, that the Tort Immunity Act was applicable, and
%that plaintiff had alleged insufficient facts to establish a special duty owed
%to plaintiff's insured. The trial court granted defendant's motion and
%dismissed plaintiff's complaint with prejudice.

We\ldots
%initially
consider whether plaintiff's complaint adequately set forth a claim
under a bailment theory.

A bailment is the delivery of property for some purpose upon a contract, express
or implied, that after the purpose has been fulfilled, the property shall be
redelivered to the bailor, or otherwise dealt with according to his directions,
or kept until he reclaims it. \textit{Kirby v. Chicago City Bank \& Trust Co.},
82 Ill. App. 3d 1113, 1116 (1980).
In order to properly plead the existence of and the right to recover under a
bailment theory, the following elements must be alleged: an agreement, express
or implied, to create a bailment; delivery of the property in good condition;
acceptance of the property bailed by the bailee; and nonreturn or redelivery of
the property in a damaged condition.\ldots
%\textit{Mueller v. Soffer} (1987), 160 Ill.
%App.3d 699, 704, 112 Ill.Dec. 589, 592, 513 N.E.2d 1198, 1201; \textit{Robinson
%v. St. Clair County} (1986), 144 Ill.App.3d 118, 120, 98 Ill.Dec. 133, 134, 493
%N.E.2d 1154, 1155; \textit{Wright v. Autohaus Fortense, Inc.} (1984), 129
%Ill.App.3d 422, 424-25, 84 Ill.Dec. 633, 635, 472 N.E.2d 593, 595.

The complaint alleged that defendant impounded the subject automobile after the
driver was arrested on December 22, 1988, and although [plaintiff]
demanded the return of his vehicle, defendant was unable or unwilling to return
it. A police officer ultimately reported the automobile stolen on January 26,
1989. Defendant did not file an answer denying these assertions, and its motion
to dismiss admitted all well pleaded facts.
%\textit{Petrauskas v. Wexenthaller
%Realty Management, Inc.} (1989), 186 Ill.App.3d 820, 825, 134 Ill.Dec. 556,
%559, 542 N.E.2d 902, 905.

On this record, it is undisputed that [plaintiff's] automobile
%of the plaintiff's insured
was actually transferred to and accepted by defendant when the vehicle was
impounded by the police department at the time of the driver's arrest. It is
also undisputed that the vehicle was never returned to the [plaintiff]
%'s insured
and the defendant exercised exclusive control over the vehicle from the time it
was impounded until it was stolen from the police department's lot.

Defendant argues that plaintiff cannot recover under a bailment theory because
it has not established an agreement by its insured and the police department to
create a bailment. It has been held, however, that a bailment is a consensual
relationship that can be established by express contract or by implication,
\textit{Kirby}, 82 Ill. App. 3d at 1116,
%38 Ill.Dec. at 492, 403 N.E.2d at 723;
%textit{Berglund v. Roosevelt University} (1974), 18 Ill.App.3d 842, 844, 310
%N.E.2d 773, 776),
and the agreement of the parties to a bailment may be implied
in law or in fact, \textit{Robinson
v. St. Clair County}, 144 Ill. App. 3d 118, 120 (1986).
%  (\textit{Robinson}, 144 Ill.App.3d at 120, 98 Ill.Dec. at 134,
%493 N.E.2d at 1155; \textit{Berglund}, 18 Ill.App.3d at 844, 310 N.E.2d  at
%776; \textit{Chesterfield Sewer \& Water, Inc. v. Citizens Insurance Company}
%(1965), 57 Ill.App.2d 90, 94-95, 207 N.E.2d 84, 86).

A constructive bailment, or a bailment implied in law, may be found where the
property of one person is voluntarily received by another for some purpose other
than that of obtaining ownership.
%(\textit{Chesterfield}, 57 Ill.App.2d at 93, 207 N.E.2d at 86.)
In such cases, the law implies a contract for the keeping of
the property until it shall be restored to the owner or his agent, and the
contract implied is that of a depositary. Thus, the duty of a depositary is
imposed without any actual contract for that purpose.
%(\textit{Chesterfield}, 57 Ill.App.2d at 93, 207 N.E.2d at 86.)
The holder is bound to take care of, keep,
and preserve the property, not for the sake of any benefit to himself or upon
any expectation of compensation for his services, but solely for the convenience
and accommodation of the owner.\ldots
%\textit{Chesterfield}, 57 Ill.App.2d at 93, 207 N.E.2d at 86.

%The court noted in \textit{Chesterfield} that
%
%\begin{quote} ``[w]here, otherwise than by mutual contract of bailment, one
%person has lawfully acquired the possession of personal property of another and
%holds it under circumstances whereby he ought, upon principles of justice, to
%keep it safely and restore it or deliver it to the owner, such person, and the
%owner of the property are, by operation of law, generally treated as bailee and
%bailor under a contract of bailment, irrespective of whether or not there has
%been any mutual assent, expressed or implied, to such relationship.''
%\textit{Chesterfield}, 57 Ill.App.2d at 94, 207 N.E.2d at 86, quoting
%\textit{Woodson v. Hare} (1943), 244 Ala. 301, 13 So.2d 172, 174.
%\end{quote}

We
find that the assertions in plaintiff's complaint were sufficient to allege a
constructive bailment. The complaint asserted that the automobile of the
plaintiff's insured was impounded when the driver was arrested on December 22,
1988. Although the insured demanded the return of the car, the police department
was unable or unwilling to return it. The police department exercised exclusive
control over the vehicle until it was stolen. Thus, the police department
lawfully acquired possession of the vehicle and held it under circumstances
whereby it was obligated, under principles of justice, to keep it safely and
restore it or deliver it to the owner.

%See \textit{Chesterfield}, 57 Ill.App.2d
%at 93-94, 207 N.E.2d at 86; \textit{City of St. Paul v. Myles} (1974), 298
%Minn. 298, 300, 218 N.W.2d 697, 699.

%We next address defendant's assertion that the Tort Immunity Act was applicable
%and precluded recovery by the plaintiff.
%
%Defendant's motion to dismiss the complaint was predicated upon section 4-102 of
%the Tort Immunity Act which provides in pertinent part as follows:
%
%\begin{quote} ``Neither a local public entity nor a public employee is liable
%for failure to establish a police department or otherwise provide police
%protection service or, if police protection service is provided, for failure to
%provide adequate police protection or service, failure to prevent the commission
%of crimes, failure to detect or solve crimes, and failure to identify or
%apprehend criminals.'' Ill.Rev. Stat.1987, ch. 85, par. 4-102. \end{quote}
%Contrary to the assertions in defendant's motion to dismiss, plaintiff's claim
%was not based upon the defendant's failure to provide police protection or the
%failure to detect or prevent the commission of a crime. Rather, the complaint
%was premised upon bailment, a contractual relationship, and asserted that
%defendant accepted, exercised exclusive possession of, but failed to return the
%automobile of the plaintiff's insured.
%
%The Tort Immunity Act specifically states that it provides local public entities
%and public employees immunity from liability arising from the operation of
%government (Ill.Rev.Stat.1987, ch. 85, par. 1-101.1) but does not in any way
%affect the liability of a local public entity or public employee based on
%contract (Ill.Rev.Stat.1987, ch. 85, par. 2-101). Thus, the Act applies only to
%tort actions. \textit{Firestone v. Fritz} (1983), 119 Ill.App.3d 685, 689, 75
%Ill.Dec. 83, 87, 456 N.E.2d 904, 908.
%
% In construing statutes, the role of the courts is to ascertain and give effect
%to the intent of the legislature. (\textit{Methodist Medical Center of Illinois
%v. Taylor} (1986), 140 Ill.App.3d 713, 716, 95 Ill.Dec. 130, 133, 489 N.E.2d
%351, 354; \textit{Indian Valley Golf Club, Inc. v. Village of Long Grove}
%(1985), 135 Ill.App.3d 543, 89 Ill.Dec. 726, 481 N.E.2d 277.) In doing so,
%however, the court should not create new rights or limitations not suggested by
%the statutory language. (\textit{Methodist Medical Center of Illinois}, 140
%Ill.App.3d at 716-17, 95 Ill.Dec. at 133, 489 N.E.2d at 354; \textit{Jackson v.
%Navik} (1974), 17 Ill.App.3d 672, 676, 308 N.E.2d 143, 146.) The plain meaning
%of the language used by the legislature is the safest guide in construing any
%act, and the court has no right to read into the statute words that are not
%found there either by express inclusion or by fair implication. (\textit{Indian
%Valley Golf Club, Inc.}, 135 Ill.App.3d at 552, 89 Ill.Dec. at 732, 481 N.E.2d
%at 283; \textit{Illinois Publishing \& Printing Co. v. Industrial Commission}
%(1921), 299 Ill. 189, 196, 132 N.E. 511.) Additionally, because the Tort
%Immunity Act is in derogation of the common law, the statute should be strictly
%construed against the local public entity. \textit{Reynolds v. City of Tuscola}
%(1971), 48 Ill.2d 339, 342, 270 N.E.2d 415, 417; \textit{Clark v. City of
%Chicago} (1980), 88 Ill. App.3d 760, 764, 43 Ill.Dec. 892, 896, 410 N.E.2d 1025,
%1029; \textit{Mounce v. City of Lincoln} (1978), 64 Ill.App.3d 461, 463, 21 Ill.
%Dec. 312, 313, 381 N.E.2d 421, 422.
%
%Nothing in the language of the Tort Immunity Act provides local public entities
%or public employees immunity from suit for breach of contract or the failure to
%return bailed property in good condition. To infer such immunity would be an
%abuse of the court's responsibility for statutory construction.
%(\textit{Methodist Medical Center of Illinois}, 140 Ill.App.3d at 717, 95
%Ill.Dec. at 133, 489 N.E.2d at 354.) Consequently, the Tort Immunity Act is
%inapplicable under the facts at bar, and defendant's motion to dismiss based
%upon section 4-102 of that act should have been denied.
%
%Defendant argues that plaintiff's claim must be predicated upon the tort theory
%of negligence. We find this argument unpersuasive. Plaintiff's contract action
%here, predicated upon a bailment, is not transformed into a tort action merely
%by characterizing the conduct which caused the breach as negligent.
%
%Based upon our resolution of these issues, we need not address plaintiff's claim
%that defendant was liable for breach of a special duty owed to its insured.

%For the foregoing reasons, the order of the circuit court of Cook County is
%reversed, and the cause is remanded for further proceedings consistent with the
%views expressed herein.

